% =====================================================================
%  BUỔI 4 — LỜI GIẢI BÀI TẬP VỀ NHÀ (từ giao-vien/02-loi-giai-bai-tap-ve-nha.md)
%  Sinh tự động bằng scripts/md2tex-luoigiai.py — không sửa tay.
%  Tài liệu giáo viên: dùng để chữa bài trên lớp và soạn đề kiểm tra.
% =====================================================================
\documentclass[11pt, a4paper]{article}
% =====================================================================
%  CAE SHANGHAI — GÓI LỆNH DÙNG CHUNG CHO TÀI LIỆU LUYỆN THI CSCA TOÁN
%  Biên dịch bằng XeLaTeX. Mọi file .tex trong buổi học \input file này.
% =====================================================================

% ---------- NGÔN NGỮ & PHÔNG CHỮ ----------
\usepackage{fontspec}
\usepackage{polyglossia}
\setmainlanguage{vietnamese}
\setotherlanguage{english}

% Chữ Latin (Anh + Việt): TeX Gyre Termes — bản clone Times New Roman,
% đi kèm TeX Live, hỗ trợ đầy đủ dấu tiếng Việt.
\setmainfont{TeX Gyre Termes}
\setsansfont{TeX Gyre Heros}

% Chữ Hán: ưu tiên Source Han Serif SC (思源宋体) — chuẩn học thuật Trung Quốc.
% Nếu máy chưa cài, tự động dùng Songti SC (宋体, có sẵn trên macOS).
% Cả hai đều thuộc họ 宋体 nên diện mạo tài liệu không đổi.
\IfFontExistsTF{Source Han Serif SC}{
  \newcommand{\cjkfamilyname}{Source Han Serif SC}
}{
  \IfFontExistsTF{Songti SC}{
    \newcommand{\cjkfamilyname}{Songti SC}
  }{
    \newcommand{\cjkfamilyname}{Noto Serif CJK SC}
  }
}

% xeCJK tự chuyển font khi gặp ký tự Hán, không cần bọc \cjk{} thủ công.
% Đây là điều kiện bắt buộc: font Latin (TeX Gyre Termes) không có glyph chữ Hán,
% nếu không khai báo xeCJK thì mọi chữ Hán viết trần sẽ bị nuốt khỏi bản in.
\usepackage{xeCJK}
\setCJKfamilyfont{caehan}{\cjkfamilyname}[Scale=0.95]
\setCJKmainfont{\cjkfamilyname}[Scale=0.95]
\setCJKsansfont{\cjkfamilyname}[Scale=0.95]
\setCJKmonofont{\cjkfamilyname}[Scale=0.95]

% \cjk{} ép dùng font Hán trong ngữ cảnh cần thiết (ví dụ trong tiêu đề in đậm
% hoặc trong bảng có định dạng riêng). Dùng \CJKfamily chứ không phải tên font,
% vì tên font chỉ là văn bản — sẽ in ra chữ thay vì đổi font.
\newcommand{\cjk}[1]{{\CJKfamily{caehan}#1}}

% ---------- TOÁN HỌC ----------
\usepackage{amsmath, amssymb, amsthm, mathtools}
\usepackage{unicode-math}
\setmathfont{TeX Gyre Termes Math}

% ---------- BỐ CỤC TRANG ----------
% Lề phải chừa vùng an toàn của letterhead:
%   banner trên chiếm 0–19mm  → top tối thiểu 2,2cm
%   banner dưới chiếm 281–297mm → bottom tối thiểu 2,2cm
% Lề trên 2,6cm (26mm) để tiêu đề đầu trang không đè banner (banner hết ở 19mm).
\usepackage[a4paper, top=2.6cm, bottom=2.4cm, left=2.5cm, right=2.0cm]{geometry}
\usepackage{fancyhdr}
\usepackage{titlesec}
\usepackage{enumitem}
\usepackage{graphicx}
\usepackage{booktabs, array, longtable, tabularx}
\usepackage{xcolor}
\usepackage[most]{tcolorbox}
\usepackage[colorlinks=true, linkcolor=cae_blue, citecolor=cae_blue, urlcolor=cae_blue]{hyperref}

% ---------- BẢNG MÀU THƯƠNG HIỆU ----------
\definecolor{cae_blue}{RGB}{0, 51, 102}
\definecolor{cae_gold}{RGB}{196, 149, 60}
\definecolor{box_bg}{RGB}{245, 247, 250}
\definecolor{box_line}{RGB}{0, 51, 102}

% ---------- TIÊU ĐỀ MỤC ----------
% Đệm đầu trang: banner letterhead chiếm 0–19mm, mà lề trên là 22mm.
% \section* đầu trang (sau \clearpage) sát lề nên phải chừa thêm, nếu không
% sẽ đè lên banner.
\titleformat{\section}
  {\normalfont\large\bfseries\sffamily\color{cae_blue}}{\thesection.}{0.6em}{}
\titleformat{\subsection}
  {\normalfont\normalsize\bfseries\sffamily\color{cae_blue!85!black}}{\thesubsection}{0.6em}{}
\titlespacing*{\section}{0pt}{14pt}{6pt}
\titlespacing*{\subsection}{0pt}{10pt}{4pt}

% Tiêu đề không đánh số ở đầu trang: \section* sau \clearpage nằm sát lề trên
% nên đè lên banner letterhead. Dùng \vspace* ngay TRƯỚC \section* không đủ,
% vì titlesec chèn khoảng cách riêng. Cách chắc chắn là đệm bằng \null.
\newcommand{\sectiondautrang}[1]{\section*{#1}}

% Dạng bài dùng \subsection*{} nên cần định dạng riêng, đậm và có gạch chân màu
\titleformat{\subsubsection}
  {\normalfont\normalsize\bfseries\sffamily\color{cae_blue!85!black}}{\thesubsubsection}{0.6em}{}

% ---------- LETTERHEAD (ảnh nền toàn trang) ----------
% Ảnh letterhead của CAE Shanghai: banner xanh trên/dưới kèm logo chìm giữa.
% Tỉ lệ ảnh 1414x2000 = 0.7070, khớp A4 (0.7071) nên chèn thẳng không méo.
% Đặt ở lớp nền dưới cùng để chữ không bị ảnh che.
% Dùng \AddToShipoutPictureBG (không có dấu *) để áp dụng cho MỌI trang.
\usepackage{eso-pic}
\AddToShipoutPictureBG{%
  \AtPageLowerLeft{%
    \includegraphics[width=\paperwidth,height=\paperheight]{images/letterhead.png}%
  }%
}

% ---------- ĐẦU / CHÂN TRANG ----------
% Letterhead đã có banner trên và dưới, nên bỏ dòng chạy để tránh chồng chữ.
% Số trang đặt trong vùng nội dung, sát trên banner dưới, lề phải.
\pagestyle{fancy}
\setlength{\headheight}{0pt}
\fancyhf{}
\renewcommand{\headrulewidth}{0pt}
\renewcommand{\footrulewidth}{0pt}
\fancyfoot[R]{\raisebox{0.9cm}{\small\color{cae_blue}\thepage}}

% ---------- MÔI TRƯỜNG NỘI DUNG ----------
% Trong tcolorbox, LaTeX đặt parindent = parskip = 0, nên các đoạn văn trong
% khung dính liền nhau. Đặt lại khoảng cách đoạn cho mọi khung.
\newcommand{\khoangcachdoan}{\setlength{\parskip}{0.5em}\setlength{\parindent}{0pt}}
\tcbset{before upper={\khoangcachdoan}}

\newtcolorbox{dinhnghia}[1][]{
  colback=box_bg, colframe=box_line, boxrule=0.6pt, arc=1pt,
  fonttitle=\bfseries, title={Định nghĩa #1}, breakable}
\newtcolorbox{dinhly}[1][]{
  colback=box_bg, colframe=cae_gold, boxrule=0.6pt, arc=1pt,
  fonttitle=\bfseries, title={Tính chất #1}, breakable}
\newtcolorbox{vidu}[1][]{
  colback=white, colframe=black!40, boxrule=0.4pt, arc=1pt,
  fonttitle=\bfseries, title={Ví dụ #1}, breakable}
\newtcolorbox{baitap}[1][]{
  colback=white, colframe=cae_blue!50, boxrule=0.4pt, arc=1pt,
  fonttitle=\bfseries, title={Câu #1}, breakable}

% Ô đáp án: viền xanh đậm hơn, dùng trong file đáp án
\newtcolorbox{dapan}[1][]{
  colback=box_bg, colframe=cae_blue, boxrule=0.6pt, arc=1pt,
  fonttitle=\bfseries, title={Câu #1}, breakable}

% ---------- TÊN CHỨNG MINH ----------
\renewcommand{\proofname}{\normalfont\bfseries Chứng minh}

% ---------- MÔN HỌC ----------
% Đổi được ở file .tex trước khi % =====================================================================
%  CAE SHANGHAI — GÓI LỆNH DÙNG CHUNG CHO TÀI LIỆU LUYỆN THI CSCA TOÁN
%  Biên dịch bằng XeLaTeX. Mọi file .tex trong buổi học \input file này.
% =====================================================================

% ---------- NGÔN NGỮ & PHÔNG CHỮ ----------
\usepackage{fontspec}
\usepackage{polyglossia}
\setmainlanguage{vietnamese}
\setotherlanguage{english}

% Chữ Latin (Anh + Việt): TeX Gyre Termes — bản clone Times New Roman,
% đi kèm TeX Live, hỗ trợ đầy đủ dấu tiếng Việt.
\setmainfont{TeX Gyre Termes}
\setsansfont{TeX Gyre Heros}

% Chữ Hán: ưu tiên Source Han Serif SC (思源宋体) — chuẩn học thuật Trung Quốc.
% Nếu máy chưa cài, tự động dùng Songti SC (宋体, có sẵn trên macOS).
% Cả hai đều thuộc họ 宋体 nên diện mạo tài liệu không đổi.
\IfFontExistsTF{Source Han Serif SC}{
  \newcommand{\cjkfamilyname}{Source Han Serif SC}
}{
  \IfFontExistsTF{Songti SC}{
    \newcommand{\cjkfamilyname}{Songti SC}
  }{
    \newcommand{\cjkfamilyname}{Noto Serif CJK SC}
  }
}

% xeCJK tự chuyển font khi gặp ký tự Hán, không cần bọc \cjk{} thủ công.
% Đây là điều kiện bắt buộc: font Latin (TeX Gyre Termes) không có glyph chữ Hán,
% nếu không khai báo xeCJK thì mọi chữ Hán viết trần sẽ bị nuốt khỏi bản in.
\usepackage{xeCJK}
\setCJKfamilyfont{caehan}{\cjkfamilyname}[Scale=0.95]
\setCJKmainfont{\cjkfamilyname}[Scale=0.95]
\setCJKsansfont{\cjkfamilyname}[Scale=0.95]
\setCJKmonofont{\cjkfamilyname}[Scale=0.95]

% \cjk{} ép dùng font Hán trong ngữ cảnh cần thiết (ví dụ trong tiêu đề in đậm
% hoặc trong bảng có định dạng riêng). Dùng \CJKfamily chứ không phải tên font,
% vì tên font chỉ là văn bản — sẽ in ra chữ thay vì đổi font.
\newcommand{\cjk}[1]{{\CJKfamily{caehan}#1}}

% ---------- TOÁN HỌC ----------
\usepackage{amsmath, amssymb, amsthm, mathtools}
\usepackage{unicode-math}
\setmathfont{TeX Gyre Termes Math}

% ---------- BỐ CỤC TRANG ----------
% Lề phải chừa vùng an toàn của letterhead:
%   banner trên chiếm 0–19mm  → top tối thiểu 2,2cm
%   banner dưới chiếm 281–297mm → bottom tối thiểu 2,2cm
% Lề trên 2,6cm (26mm) để tiêu đề đầu trang không đè banner (banner hết ở 19mm).
\usepackage[a4paper, top=2.6cm, bottom=2.4cm, left=2.5cm, right=2.0cm]{geometry}
\usepackage{fancyhdr}
\usepackage{titlesec}
\usepackage{enumitem}
\usepackage{graphicx}
\usepackage{booktabs, array, longtable, tabularx}
\usepackage{xcolor}
\usepackage[most]{tcolorbox}
\usepackage[colorlinks=true, linkcolor=cae_blue, citecolor=cae_blue, urlcolor=cae_blue]{hyperref}

% ---------- BẢNG MÀU THƯƠNG HIỆU ----------
\definecolor{cae_blue}{RGB}{0, 51, 102}
\definecolor{cae_gold}{RGB}{196, 149, 60}
\definecolor{box_bg}{RGB}{245, 247, 250}
\definecolor{box_line}{RGB}{0, 51, 102}

% ---------- TIÊU ĐỀ MỤC ----------
% Đệm đầu trang: banner letterhead chiếm 0–19mm, mà lề trên là 22mm.
% \section* đầu trang (sau \clearpage) sát lề nên phải chừa thêm, nếu không
% sẽ đè lên banner.
\titleformat{\section}
  {\normalfont\large\bfseries\sffamily\color{cae_blue}}{\thesection.}{0.6em}{}
\titleformat{\subsection}
  {\normalfont\normalsize\bfseries\sffamily\color{cae_blue!85!black}}{\thesubsection}{0.6em}{}
\titlespacing*{\section}{0pt}{14pt}{6pt}
\titlespacing*{\subsection}{0pt}{10pt}{4pt}

% Tiêu đề không đánh số ở đầu trang: \section* sau \clearpage nằm sát lề trên
% nên đè lên banner letterhead. Dùng \vspace* ngay TRƯỚC \section* không đủ,
% vì titlesec chèn khoảng cách riêng. Cách chắc chắn là đệm bằng \null.
\newcommand{\sectiondautrang}[1]{\section*{#1}}

% Dạng bài dùng \subsection*{} nên cần định dạng riêng, đậm và có gạch chân màu
\titleformat{\subsubsection}
  {\normalfont\normalsize\bfseries\sffamily\color{cae_blue!85!black}}{\thesubsubsection}{0.6em}{}

% ---------- LETTERHEAD (ảnh nền toàn trang) ----------
% Ảnh letterhead của CAE Shanghai: banner xanh trên/dưới kèm logo chìm giữa.
% Tỉ lệ ảnh 1414x2000 = 0.7070, khớp A4 (0.7071) nên chèn thẳng không méo.
% Đặt ở lớp nền dưới cùng để chữ không bị ảnh che.
% Dùng \AddToShipoutPictureBG (không có dấu *) để áp dụng cho MỌI trang.
\usepackage{eso-pic}
\AddToShipoutPictureBG{%
  \AtPageLowerLeft{%
    \includegraphics[width=\paperwidth,height=\paperheight]{images/letterhead.png}%
  }%
}

% ---------- ĐẦU / CHÂN TRANG ----------
% Letterhead đã có banner trên và dưới, nên bỏ dòng chạy để tránh chồng chữ.
% Số trang đặt trong vùng nội dung, sát trên banner dưới, lề phải.
\pagestyle{fancy}
\setlength{\headheight}{0pt}
\fancyhf{}
\renewcommand{\headrulewidth}{0pt}
\renewcommand{\footrulewidth}{0pt}
\fancyfoot[R]{\raisebox{0.9cm}{\small\color{cae_blue}\thepage}}

% ---------- MÔI TRƯỜNG NỘI DUNG ----------
% Trong tcolorbox, LaTeX đặt parindent = parskip = 0, nên các đoạn văn trong
% khung dính liền nhau. Đặt lại khoảng cách đoạn cho mọi khung.
\newcommand{\khoangcachdoan}{\setlength{\parskip}{0.5em}\setlength{\parindent}{0pt}}
\tcbset{before upper={\khoangcachdoan}}

\newtcolorbox{dinhnghia}[1][]{
  colback=box_bg, colframe=box_line, boxrule=0.6pt, arc=1pt,
  fonttitle=\bfseries, title={Định nghĩa #1}, breakable}
\newtcolorbox{dinhly}[1][]{
  colback=box_bg, colframe=cae_gold, boxrule=0.6pt, arc=1pt,
  fonttitle=\bfseries, title={Tính chất #1}, breakable}
\newtcolorbox{vidu}[1][]{
  colback=white, colframe=black!40, boxrule=0.4pt, arc=1pt,
  fonttitle=\bfseries, title={Ví dụ #1}, breakable}
\newtcolorbox{baitap}[1][]{
  colback=white, colframe=cae_blue!50, boxrule=0.4pt, arc=1pt,
  fonttitle=\bfseries, title={Câu #1}, breakable}

% Ô đáp án: viền xanh đậm hơn, dùng trong file đáp án
\newtcolorbox{dapan}[1][]{
  colback=box_bg, colframe=cae_blue, boxrule=0.6pt, arc=1pt,
  fonttitle=\bfseries, title={Câu #1}, breakable}

% ---------- TÊN CHỨNG MINH ----------
\renewcommand{\proofname}{\normalfont\bfseries Chứng minh}

% ---------- MÔN HỌC ----------
% Đổi được ở file .tex trước khi % =====================================================================
%  CAE SHANGHAI — GÓI LỆNH DÙNG CHUNG CHO TÀI LIỆU LUYỆN THI CSCA TOÁN
%  Biên dịch bằng XeLaTeX. Mọi file .tex trong buổi học \input file này.
% =====================================================================

% ---------- NGÔN NGỮ & PHÔNG CHỮ ----------
\usepackage{fontspec}
\usepackage{polyglossia}
\setmainlanguage{vietnamese}
\setotherlanguage{english}

% Chữ Latin (Anh + Việt): TeX Gyre Termes — bản clone Times New Roman,
% đi kèm TeX Live, hỗ trợ đầy đủ dấu tiếng Việt.
\setmainfont{TeX Gyre Termes}
\setsansfont{TeX Gyre Heros}

% Chữ Hán: ưu tiên Source Han Serif SC (思源宋体) — chuẩn học thuật Trung Quốc.
% Nếu máy chưa cài, tự động dùng Songti SC (宋体, có sẵn trên macOS).
% Cả hai đều thuộc họ 宋体 nên diện mạo tài liệu không đổi.
\IfFontExistsTF{Source Han Serif SC}{
  \newcommand{\cjkfamilyname}{Source Han Serif SC}
}{
  \IfFontExistsTF{Songti SC}{
    \newcommand{\cjkfamilyname}{Songti SC}
  }{
    \newcommand{\cjkfamilyname}{Noto Serif CJK SC}
  }
}

% xeCJK tự chuyển font khi gặp ký tự Hán, không cần bọc \cjk{} thủ công.
% Đây là điều kiện bắt buộc: font Latin (TeX Gyre Termes) không có glyph chữ Hán,
% nếu không khai báo xeCJK thì mọi chữ Hán viết trần sẽ bị nuốt khỏi bản in.
\usepackage{xeCJK}
\setCJKfamilyfont{caehan}{\cjkfamilyname}[Scale=0.95]
\setCJKmainfont{\cjkfamilyname}[Scale=0.95]
\setCJKsansfont{\cjkfamilyname}[Scale=0.95]
\setCJKmonofont{\cjkfamilyname}[Scale=0.95]

% \cjk{} ép dùng font Hán trong ngữ cảnh cần thiết (ví dụ trong tiêu đề in đậm
% hoặc trong bảng có định dạng riêng). Dùng \CJKfamily chứ không phải tên font,
% vì tên font chỉ là văn bản — sẽ in ra chữ thay vì đổi font.
\newcommand{\cjk}[1]{{\CJKfamily{caehan}#1}}

% ---------- TOÁN HỌC ----------
\usepackage{amsmath, amssymb, amsthm, mathtools}
\usepackage{unicode-math}
\setmathfont{TeX Gyre Termes Math}

% ---------- BỐ CỤC TRANG ----------
% Lề phải chừa vùng an toàn của letterhead:
%   banner trên chiếm 0–19mm  → top tối thiểu 2,2cm
%   banner dưới chiếm 281–297mm → bottom tối thiểu 2,2cm
% Lề trên 2,6cm (26mm) để tiêu đề đầu trang không đè banner (banner hết ở 19mm).
\usepackage[a4paper, top=2.6cm, bottom=2.4cm, left=2.5cm, right=2.0cm]{geometry}
\usepackage{fancyhdr}
\usepackage{titlesec}
\usepackage{enumitem}
\usepackage{graphicx}
\usepackage{booktabs, array, longtable, tabularx}
\usepackage{xcolor}
\usepackage[most]{tcolorbox}
\usepackage[colorlinks=true, linkcolor=cae_blue, citecolor=cae_blue, urlcolor=cae_blue]{hyperref}

% ---------- BẢNG MÀU THƯƠNG HIỆU ----------
\definecolor{cae_blue}{RGB}{0, 51, 102}
\definecolor{cae_gold}{RGB}{196, 149, 60}
\definecolor{box_bg}{RGB}{245, 247, 250}
\definecolor{box_line}{RGB}{0, 51, 102}

% ---------- TIÊU ĐỀ MỤC ----------
% Đệm đầu trang: banner letterhead chiếm 0–19mm, mà lề trên là 22mm.
% \section* đầu trang (sau \clearpage) sát lề nên phải chừa thêm, nếu không
% sẽ đè lên banner.
\titleformat{\section}
  {\normalfont\large\bfseries\sffamily\color{cae_blue}}{\thesection.}{0.6em}{}
\titleformat{\subsection}
  {\normalfont\normalsize\bfseries\sffamily\color{cae_blue!85!black}}{\thesubsection}{0.6em}{}
\titlespacing*{\section}{0pt}{14pt}{6pt}
\titlespacing*{\subsection}{0pt}{10pt}{4pt}

% Tiêu đề không đánh số ở đầu trang: \section* sau \clearpage nằm sát lề trên
% nên đè lên banner letterhead. Dùng \vspace* ngay TRƯỚC \section* không đủ,
% vì titlesec chèn khoảng cách riêng. Cách chắc chắn là đệm bằng \null.
\newcommand{\sectiondautrang}[1]{\section*{#1}}

% Dạng bài dùng \subsection*{} nên cần định dạng riêng, đậm và có gạch chân màu
\titleformat{\subsubsection}
  {\normalfont\normalsize\bfseries\sffamily\color{cae_blue!85!black}}{\thesubsubsection}{0.6em}{}

% ---------- LETTERHEAD (ảnh nền toàn trang) ----------
% Ảnh letterhead của CAE Shanghai: banner xanh trên/dưới kèm logo chìm giữa.
% Tỉ lệ ảnh 1414x2000 = 0.7070, khớp A4 (0.7071) nên chèn thẳng không méo.
% Đặt ở lớp nền dưới cùng để chữ không bị ảnh che.
% Dùng \AddToShipoutPictureBG (không có dấu *) để áp dụng cho MỌI trang.
\usepackage{eso-pic}
\AddToShipoutPictureBG{%
  \AtPageLowerLeft{%
    \includegraphics[width=\paperwidth,height=\paperheight]{images/letterhead.png}%
  }%
}

% ---------- ĐẦU / CHÂN TRANG ----------
% Letterhead đã có banner trên và dưới, nên bỏ dòng chạy để tránh chồng chữ.
% Số trang đặt trong vùng nội dung, sát trên banner dưới, lề phải.
\pagestyle{fancy}
\setlength{\headheight}{0pt}
\fancyhf{}
\renewcommand{\headrulewidth}{0pt}
\renewcommand{\footrulewidth}{0pt}
\fancyfoot[R]{\raisebox{0.9cm}{\small\color{cae_blue}\thepage}}

% ---------- MÔI TRƯỜNG NỘI DUNG ----------
% Trong tcolorbox, LaTeX đặt parindent = parskip = 0, nên các đoạn văn trong
% khung dính liền nhau. Đặt lại khoảng cách đoạn cho mọi khung.
\newcommand{\khoangcachdoan}{\setlength{\parskip}{0.5em}\setlength{\parindent}{0pt}}
\tcbset{before upper={\khoangcachdoan}}

\newtcolorbox{dinhnghia}[1][]{
  colback=box_bg, colframe=box_line, boxrule=0.6pt, arc=1pt,
  fonttitle=\bfseries, title={Định nghĩa #1}, breakable}
\newtcolorbox{dinhly}[1][]{
  colback=box_bg, colframe=cae_gold, boxrule=0.6pt, arc=1pt,
  fonttitle=\bfseries, title={Tính chất #1}, breakable}
\newtcolorbox{vidu}[1][]{
  colback=white, colframe=black!40, boxrule=0.4pt, arc=1pt,
  fonttitle=\bfseries, title={Ví dụ #1}, breakable}
\newtcolorbox{baitap}[1][]{
  colback=white, colframe=cae_blue!50, boxrule=0.4pt, arc=1pt,
  fonttitle=\bfseries, title={Câu #1}, breakable}

% Ô đáp án: viền xanh đậm hơn, dùng trong file đáp án
\newtcolorbox{dapan}[1][]{
  colback=box_bg, colframe=cae_blue, boxrule=0.6pt, arc=1pt,
  fonttitle=\bfseries, title={Câu #1}, breakable}

% ---------- TÊN CHỨNG MINH ----------
\renewcommand{\proofname}{\normalfont\bfseries Chứng minh}

% ---------- MÔN HỌC ----------
% Đổi được ở file .tex trước khi \input{cae-style} bằng \renewcommand{\monhoc}{...}
\providecommand{\monhoc}{Môn Toán}

% ---------- TRANG BÌA GỌN ----------
% #1 = loại tài liệu (TÀI LIỆU CHUẨN BỊ, TÀI LIỆU BUỔI HỌC, ...)
% #2 = số buổi
% #3 = tên buổi tiếng Việt
% #4 = tên buổi tiếng Hán
% #5 = module và tỉ trọng
\newcommand{\trangbia}[5]{%
  \begin{center}
    \vspace*{0.6cm}
    {\large\bfseries\color{cae_blue} CAE SHANGHAI}\\[0.15cm]
    {\footnotesize Tài liệu luyện thi CSCA · \monhoc}\\[0.5cm]
    {\color{cae_gold}\rule{0.55\linewidth}{0.8pt}}\\[0.35cm]
    {\small\bfseries\color{cae_blue!80!black} #1}\\[0.35cm]
    {\Large\bfseries\color{cae_blue} BUỔI #2 — #3}\\[0.2cm]
    {\large\color{cae_blue!85!black} \cjk{#4}}\\[0.25cm]
    \ifx\relax#5\relax\else{\small #5}\\[0.35cm]\fi
    {\color{cae_gold}\rule{0.55\linewidth}{0.8pt}}\\[0.3cm]
    {\footnotesize Biên soạn: Hoàng Dương, CAE Shanghai \quad·\quad \today}
  \end{center}
  \vspace{0.3cm}
}

% ---------- TRANG BÌA CHO ĐỀ KIỂM TRA ----------
% Khác \trangbia ở chỗ không có chữ "BUỔI N" — dùng cho đề thi, không phải bài giảng.
% #1 = loại tài liệu, #2 = tên kỳ thi (Việt), #3 = tên kỳ thi (Hán), #4 = thông số
\newcommand{\biade}[4]{%
  \begin{center}
    \vspace*{0.6cm}
    {\large\bfseries\color{cae_blue} CAE SHANGHAI}\\[0.15cm]
    {\footnotesize Tài liệu luyện thi CSCA · \monhoc}\\[0.5cm]
    {\color{cae_gold}\rule{0.55\linewidth}{0.8pt}}\\[0.35cm]
    {\small\bfseries\color{cae_blue!80!black} #1}\\[0.35cm]
    {\Large\bfseries\color{cae_blue} #2}\\[0.2cm]
    {\large\color{cae_blue!85!black} \cjk{#3}}\\[0.25cm]
    {\small #4}\\[0.35cm]
    {\color{cae_gold}\rule{0.55\linewidth}{0.8pt}}\\[0.3cm]
    {\footnotesize Biên soạn: Hoàng Dương, CAE Shanghai \quad·\quad \today}
  \end{center}
  \vspace{0.3cm}
}

% ---------- TIỆN ÍCH ----------
% Nhãn từ khóa in đậm ở đầu đoạn, dùng cho Từ vựng / Từ khóa nhận dạng / Chú ý
% Khoảng cách đoạn do khung tự đặt (xem \khoangcachdoan ở trên),
% ở đây chỉ cần tách đoạn và bỏ thụt đầu dòng.
\newcommand{\nhan}[1]{\par\noindent\textbf{#1.}\ \ignorespaces}

% Bảng gọn, không đánh số
\newcommand{\bang}[1]{%
  \begin{center}#1\end{center}}

% Đánh dấu ô trống cho học sinh điền
\newcommand{\chotrong}{\makebox[3.2cm]{\dotfill}}
 bằng \renewcommand{\monhoc}{...}
\providecommand{\monhoc}{Môn Toán}

% ---------- TRANG BÌA GỌN ----------
% #1 = loại tài liệu (TÀI LIỆU CHUẨN BỊ, TÀI LIỆU BUỔI HỌC, ...)
% #2 = số buổi
% #3 = tên buổi tiếng Việt
% #4 = tên buổi tiếng Hán
% #5 = module và tỉ trọng
\newcommand{\trangbia}[5]{%
  \begin{center}
    \vspace*{0.6cm}
    {\large\bfseries\color{cae_blue} CAE SHANGHAI}\\[0.15cm]
    {\footnotesize Tài liệu luyện thi CSCA · \monhoc}\\[0.5cm]
    {\color{cae_gold}\rule{0.55\linewidth}{0.8pt}}\\[0.35cm]
    {\small\bfseries\color{cae_blue!80!black} #1}\\[0.35cm]
    {\Large\bfseries\color{cae_blue} BUỔI #2 — #3}\\[0.2cm]
    {\large\color{cae_blue!85!black} \cjk{#4}}\\[0.25cm]
    \ifx\relax#5\relax\else{\small #5}\\[0.35cm]\fi
    {\color{cae_gold}\rule{0.55\linewidth}{0.8pt}}\\[0.3cm]
    {\footnotesize Biên soạn: Hoàng Dương, CAE Shanghai \quad·\quad \today}
  \end{center}
  \vspace{0.3cm}
}

% ---------- TRANG BÌA CHO ĐỀ KIỂM TRA ----------
% Khác \trangbia ở chỗ không có chữ "BUỔI N" — dùng cho đề thi, không phải bài giảng.
% #1 = loại tài liệu, #2 = tên kỳ thi (Việt), #3 = tên kỳ thi (Hán), #4 = thông số
\newcommand{\biade}[4]{%
  \begin{center}
    \vspace*{0.6cm}
    {\large\bfseries\color{cae_blue} CAE SHANGHAI}\\[0.15cm]
    {\footnotesize Tài liệu luyện thi CSCA · \monhoc}\\[0.5cm]
    {\color{cae_gold}\rule{0.55\linewidth}{0.8pt}}\\[0.35cm]
    {\small\bfseries\color{cae_blue!80!black} #1}\\[0.35cm]
    {\Large\bfseries\color{cae_blue} #2}\\[0.2cm]
    {\large\color{cae_blue!85!black} \cjk{#3}}\\[0.25cm]
    {\small #4}\\[0.35cm]
    {\color{cae_gold}\rule{0.55\linewidth}{0.8pt}}\\[0.3cm]
    {\footnotesize Biên soạn: Hoàng Dương, CAE Shanghai \quad·\quad \today}
  \end{center}
  \vspace{0.3cm}
}

% ---------- TIỆN ÍCH ----------
% Nhãn từ khóa in đậm ở đầu đoạn, dùng cho Từ vựng / Từ khóa nhận dạng / Chú ý
% Khoảng cách đoạn do khung tự đặt (xem \khoangcachdoan ở trên),
% ở đây chỉ cần tách đoạn và bỏ thụt đầu dòng.
\newcommand{\nhan}[1]{\par\noindent\textbf{#1.}\ \ignorespaces}

% Bảng gọn, không đánh số
\newcommand{\bang}[1]{%
  \begin{center}#1\end{center}}

% Đánh dấu ô trống cho học sinh điền
\newcommand{\chotrong}{\makebox[3.2cm]{\dotfill}}
 bằng \renewcommand{\monhoc}{...}
\providecommand{\monhoc}{Môn Toán}

% ---------- TRANG BÌA GỌN ----------
% #1 = loại tài liệu (TÀI LIỆU CHUẨN BỊ, TÀI LIỆU BUỔI HỌC, ...)
% #2 = số buổi
% #3 = tên buổi tiếng Việt
% #4 = tên buổi tiếng Hán
% #5 = module và tỉ trọng
\newcommand{\trangbia}[5]{%
  \begin{center}
    \vspace*{0.6cm}
    {\large\bfseries\color{cae_blue} CAE SHANGHAI}\\[0.15cm]
    {\footnotesize Tài liệu luyện thi CSCA · \monhoc}\\[0.5cm]
    {\color{cae_gold}\rule{0.55\linewidth}{0.8pt}}\\[0.35cm]
    {\small\bfseries\color{cae_blue!80!black} #1}\\[0.35cm]
    {\Large\bfseries\color{cae_blue} BUỔI #2 — #3}\\[0.2cm]
    {\large\color{cae_blue!85!black} \cjk{#4}}\\[0.25cm]
    \ifx\relax#5\relax\else{\small #5}\\[0.35cm]\fi
    {\color{cae_gold}\rule{0.55\linewidth}{0.8pt}}\\[0.3cm]
    {\footnotesize Biên soạn: Hoàng Dương, CAE Shanghai \quad·\quad \today}
  \end{center}
  \vspace{0.3cm}
}

% ---------- TRANG BÌA CHO ĐỀ KIỂM TRA ----------
% Khác \trangbia ở chỗ không có chữ "BUỔI N" — dùng cho đề thi, không phải bài giảng.
% #1 = loại tài liệu, #2 = tên kỳ thi (Việt), #3 = tên kỳ thi (Hán), #4 = thông số
\newcommand{\biade}[4]{%
  \begin{center}
    \vspace*{0.6cm}
    {\large\bfseries\color{cae_blue} CAE SHANGHAI}\\[0.15cm]
    {\footnotesize Tài liệu luyện thi CSCA · \monhoc}\\[0.5cm]
    {\color{cae_gold}\rule{0.55\linewidth}{0.8pt}}\\[0.35cm]
    {\small\bfseries\color{cae_blue!80!black} #1}\\[0.35cm]
    {\Large\bfseries\color{cae_blue} #2}\\[0.2cm]
    {\large\color{cae_blue!85!black} \cjk{#3}}\\[0.25cm]
    {\small #4}\\[0.35cm]
    {\color{cae_gold}\rule{0.55\linewidth}{0.8pt}}\\[0.3cm]
    {\footnotesize Biên soạn: Hoàng Dương, CAE Shanghai \quad·\quad \today}
  \end{center}
  \vspace{0.3cm}
}

% ---------- TIỆN ÍCH ----------
% Nhãn từ khóa in đậm ở đầu đoạn, dùng cho Từ vựng / Từ khóa nhận dạng / Chú ý
% Khoảng cách đoạn do khung tự đặt (xem \khoangcachdoan ở trên),
% ở đây chỉ cần tách đoạn và bỏ thụt đầu dòng.
\newcommand{\nhan}[1]{\par\noindent\textbf{#1.}\ \ignorespaces}

% Bảng gọn, không đánh số
\newcommand{\bang}[1]{%
  \begin{center}#1\end{center}}

% Đánh dấu ô trống cho học sinh điền
\newcommand{\chotrong}{\makebox[3.2cm]{\dotfill}}


\begin{document}

\trangbia{LỜI GIẢI BÀI TẬP VỀ NHÀ}{4}{Lượng giác: giá trị và công thức biến đổi}{三角函数·恒等变换}{Tài liệu giáo viên \quad·\quad 40 câu}

\subsection*{Nhóm A — Giá trị và công thức lượng giác}

\begin{dapan}[1 — Đáp án D]
\textbf{Giải.} Vì $\tan\alpha = 2$ nên $\cos\alpha \neq 0$. Chia cả tử và mẫu của phân thức cho $\cos\alpha$:

$$\frac{\sin\alpha - \cos\alpha}{\sin\alpha + \cos\alpha} = \frac{\dfrac{\sin\alpha}{\cos\alpha} - 1}{\dfrac{\sin\alpha}{\cos\alpha} + 1} = \frac{\tan\alpha - 1}{\tan\alpha + 1} = \frac{2-1}{2+1} = \frac{1}{3}$$

\textbf{Cách nhanh.} Không cần tìm $\sin\alpha$ và $\cos\alpha$ riêng lẻ. Chỉ một phép chia cho $\cos\alpha$ là ra đáp án.
\end{dapan}

\begin{dapan}[2 — Đáp án B]
\textbf{Giải.} Với mọi $x$, hàm $\cos$ nhận giá trị trong $[-1, 1]$:

$$-1 \le \cos\!\left(2x - \frac{\pi}{3}\right) \le 1$$

Nhân cả ba vế với $-1$ và đổi chiều bất đẳng thức:

$$-1 \le -\cos\!\left(2x - \frac{\pi}{3}\right) \le 1$$

Vậy tập giá trị là $[-1, 1]$.

\textbf{Chú ý.} Hệ số $-1$ ở trước hàm cos \textbf{không} làm thay đổi tập giá trị, chỉ đảo ngược thứ tự đạt cực đại và cực tiểu. Học sinh hay nhầm thành $[-2, 2]$ khi thấy có số $2$ trong biểu thức $2x$.
\end{dapan}

\begin{dapan}[3 — Đáp án A]
\textbf{Giải.} Từ $\sin^2\alpha + \cos^2\alpha = 1$:

$$\cos^2\alpha = 1 - \sin^2\alpha = 1 - \left(-\frac{1}{2}\right)^2 = 1 - \frac{1}{4} = \frac{3}{4}$$

Suy ra $|\cos\alpha| = \dfrac{\sqrt{3}}{2}$. Góc $\alpha$ ở phần tư thứ ba nên $\cos\alpha < 0$, vậy:

$$\cos\alpha = -\frac{\sqrt{3}}{2}$$

\textbf{Bẫy.} Học sinh quên dấu trừ và chọn đáp án B. Phần tư thứ ba có cả sin và cos đều âm.
\end{dapan}

\begin{dapan}[4 — Đáp án B]
\textbf{Giải.} Áp dụng công thức gộp bậc nhất với $a = b = 1$:

$$y = \sin x + \cos x = \sqrt{1^2+1^2}\,\sin\!\left(x + \frac{\pi}{4}\right) = \sqrt{2}\sin\!\left(x + \frac{\pi}{4}\right)$$

Vì $|\sin| \le 1$ nên giá trị lớn nhất của $y$ là $\sqrt{2}$, đạt được khi $x + \dfrac{\pi}{4} = \dfrac{\pi}{2}$, tức $x = \dfrac{\pi}{4}$.

\textbf{Cách nhanh.} Nhớ công thức $\max(a\sin x + b\cos x) = \sqrt{a^2+b^2}$, thay $a = b = 1$ là ra $\sqrt{2}$ trong vài giây.
\end{dapan}

\begin{dapan}[5 — Đáp án C]
\textbf{Giải.} Hàm cos là hàm chẵn nên $\cos(-x) = \cos x$:

$$\cos\!\left(-\frac{\pi}{6}\right) = \cos\frac{\pi}{6} = \frac{\sqrt{3}}{2}$$

\textbf{Chú ý.} Dấu trừ trong ngoặc \textbf{không} làm đổi dấu kết quả, vì cos là hàm chẵn. Chỉ với sin và tan thì dấu trừ mới đưa ra ngoài.
\end{dapan}

\begin{dapan}[6 — Đáp án D]
\textbf{Giải.} Áp dụng quy tắc đạo hàm hàm hợp: đạo hàm của $\sin u$ là $u'\cos u$ với $u = 2x + \dfrac{\pi}{4}$, và $u' = 2$.

$$f'(x) = 2\cos\!\left(2x + \frac{\pi}{4}\right)$$

\textbf{Bẫy.} Đáp án A là kết quả khi quên nhân với $u' = 2$. Đạo hàm hàm hợp luôn phải nhân thêm đạo hàm của biểu thức bên trong.
\end{dapan}

\begin{dapan}[7 — Đáp án C]
\textbf{Giải.} Thay $\sin^2 x = 1 - \cos^2 x = (1-\cos x)(1+\cos x)$:

$$\frac{\sin^2 x}{1-\cos x} = \frac{(1-\cos x)(1+\cos x)}{1-\cos x} = 1 + \cos x$$

Điều kiện $\cos x \neq 1$ đảm bảo mẫu khác $0$, nên phép rút gọn hợp lệ.

\textbf{Cách nhanh.} Thấy $\sin^2$ ở tử và một biểu thức bậc nhất của $\cos$ ở mẫu thì nghĩ ngay đến hằng đẳng thức $\sin^2 x = (1-\cos x)(1+\cos x)$.
\end{dapan}

\begin{dapan}[8 — Đáp án A]
\textbf{Giải.} Trên $\left[0, \dfrac{\pi}{2}\right]$, hàm $\sin x$ tăng từ $0$ đến $1$. Đây là một phần của nhánh đồng biến của hàm sin.
\end{dapan}

\begin{dapan}[9 — Đáp án A]
\textbf{Giải.} Hàm $y = \cos x$ là hàm chẵn: $\cos(-x) = \cos x$. Đồ thị của hàm chẵn nhận trục $Oy$ (trục $y$) làm trục đối xứng.

\textbf{Chú ý.} Đáp án B (gốc tọa độ) là đáp án dành cho hàm lẻ như $\sin x$ hay $\tan x$. Phân biệt: hàm chẵn đối xứng qua trục tung, hàm lẻ đối xứng qua gốc tọa độ.
\end{dapan}

\begin{dapan}[10 — Đáp án A]
\textbf{Giải.} Công thức nhân đôi của tan:

$$\tan 2\alpha = \frac{2\tan\alpha}{1-\tan^2\alpha}$$

\textbf{Cách nhớ.} Suy ra từ công thức cộng $\tan(\alpha+\beta)$ bằng cách thay $\beta = \alpha$, thu được $\dfrac{\tan\alpha + \tan\alpha}{1 - \tan\alpha\tan\alpha}$.

\textbf{Bẫy.} Đáp án C và D là công thức của $\cot 2\alpha$, không phải $\tan 2\alpha$. Đọc kỹ tên hàm trong đề.
\end{dapan}

\begin{dapan}[11 — Đáp án D]
\textbf{Giải.} Từ $\tan\alpha = 1$ suy ra $\sin\alpha = \cos\alpha$. Thay vào $\sin^2\alpha + \cos^2\alpha = 1$:

$$2\sin^2\alpha = 1 \implies |\sin\alpha| = \frac{\sqrt{2}}{2}$$

Góc $\alpha$ ở phần tư thứ ba nên $\sin\alpha < 0$, vậy $\sin\alpha = -\dfrac{\sqrt{2}}{2}$.

\textbf{Bẫy.} Đáp án A và B là giá trị của $\sin$ tại các góc đặc biệt khác ($\dfrac{\pi}{6}$ và $\dfrac{\pi}{3}$), không liên quan đến $\tan\alpha = 1$. Với $\tan\alpha = 1$, góc phần tư thứ ba là $\alpha = \dfrac{5\pi}{4}$.
\end{dapan}

\begin{dapan}[12 — Đáp án D]
\textbf{Giải.} Công thức cộng của sin:

$$\sin(\alpha+\beta) = \sin\alpha\cos\beta + \cos\alpha\sin\beta$$

\textbf{Chú ý.} Phân biệt bốn đáp án:
\begin{itemize}[leftmargin=1.4em, itemsep=2pt]
\item A là $\cos(\alpha+\beta) = \cos\alpha\cos\beta - \sin\alpha\sin\beta$
\item B là $\sin(\alpha-\beta) = \sin\alpha\cos\beta - \cos\alpha\sin\beta$
\item C là $\cos(\alpha-\beta) = \cos\alpha\cos\beta + \sin\alpha\sin\beta$
\item D là $\sin(\alpha+\beta)$ ✓
\end{itemize}


Học sinh hay nhầm A với D vì cả hai đều có bốn nhân tử. Quy tắc: sin thì \textbf{sin trước}, cos thì \textbf{cos trước}.
\end{dapan}

\begin{dapan}[13 — Đáp án A]
\textbf{Giải.} Hàm $y = \tan x$ đồng biến trên mỗi khoảng xác định $\left(-\dfrac{\pi}{2}+k\pi, \dfrac{\pi}{2}+k\pi\right)$.

Khoảng $\left(\dfrac{\pi}{2}, \dfrac{3\pi}{2}\right)$ là khoảng xác định ứng với $k = 1$ (dịch khoảng cơ bản sang phải một khoảng $\pi$), nên hàm tan đồng biến trên đó.

\textbf{Chú ý.} Không được nói hàm tan đồng biến trên toàn tập xác định, vì nó có các điểm gián đoạn. Chỉ đồng biến trên \textbf{từng khoảng} xác định riêng lẻ.
\end{dapan}

\subsection*{Nhóm B — Hàm bậc hai, mũ, logarit, đạo hàm}

\begin{dapan}[14 — Đáp án C]
\textbf{Giải.} Từ $f(16) = 2$ và $f(x) = \log_a x$:

$$\log\_a 16 = 2 \implies a^2 = 16 \implies a = 4 \ \text{hoặc} \ a = -4$$

Vì cơ số logarit phải dương và khác $1$, ta có $a = 4$.

\textbf{Chú ý.} Luôn loại nghiệm âm khi tìm cơ số logarit. Điều kiện $a > 0$, $a \neq 1$ là bắt buộc.
\end{dapan}

\begin{dapan}[15 — Đáp án B]
\textbf{Giải.} Từ $\log_7 x = 2$, theo định nghĩa logarit:

$$x = 7^2 = 49$$

\textbf{Bẫy.} Đáp án C ghi $7^2$ là biểu thức chưa tính, không phải giá trị cụ thể. Đáp án A ($14 = 7 \times 2$) là kết quả khi nhân sai thay vì lũy thừa.
\end{dapan}

\begin{dapan}[16 — Đáp án C]
\textbf{Giải.} Từ $f(1) = 1$ và $f(x) = \log_a(x+1)$:

$$\log\_a 2 = 1 \implies a^1 = 2 \implies a = 2$$
\end{dapan}

\begin{dapan}[17 — Đáp án B]
\textbf{Giải.} Mỗi năm giá trị còn lại $100\% - 5\% = 95\%$ của năm trước, tức nhân với hệ số $1 - 0{,}05 = 0{,}95$.

Sau $t$ năm, giá trị là:

$$V(t) = 1000 \times (0{,}95)^t = 1000(1-0{,}05)^t$$

\textbf{Bẫy.} Đáp án A là công thức \textbf{lãi đơn} (giảm đều mỗi năm một lượng như nhau). Đáp án B là công thức \textbf{lãi kép} (giảm theo phần trăm của giá trị hiện tại). Đề nói "每年折旧 5\%" nghĩa là phần trăm của giá trị hiện có, tức lãi kép.
\end{dapan}

\begin{dapan}[18 — Đáp án B]
\textbf{Giải.} Hàm $f(x) = 2^x + x - 7$ có $f'(x) = 2^x\ln 2 + 1 > 0$ với mọi $x$, nên $f$ đồng biến trên $\mathbb{R}$.

Vì $f$ đồng biến và $x_0$ là nghiệm:
\begin{itemize}[leftmargin=1.4em, itemsep=2pt]
\item Với $x_1 < x_0$ thì $f(x_1) < f(x_0) = 0$
\item Với $x_2 > x_0$ thì $f(x_2) > f(x_0) = 0$
\end{itemize}


Vậy $f(x_1) < 0$ và $f(x_2) > 0$.

\textbf{Cách nhanh.} Không cần tìm giá trị cụ thể của $x_0$. Chỉ cần biết hàm đồng biến là suy ra dấu.
\end{dapan}

\begin{dapan}[19 — Đáp án C]
\textbf{Giải.} Khai triển $f(x) = x^3 - x^2 + 5x - 5$, đạo hàm:

$$f'(x) = 3x^2 - 2x + 5$$

Thay $x = 2$:

$$f'(2) = 3(4) - 2(2) + 5 = 12 - 4 + 5 = 13$$

\textbf{Cách nhanh.} Dùng quy tắc tích: $f'(x) = 2x(x-1) + (x^2+5)$. Thay $x=2$: $4(1) + 9 = 13$.
\end{dapan}

\subsection*{Nhóm C — Tính chất hàm số}

\begin{dapan}[20 — Đáp án A]
\textbf{Giải.} Hàm $f(x) = \dfrac{ax+15}{x-a}$ là hàm lẻ khi $f(-x) = -f(x)$ với mọi $x$ trong tập xác định.

$$f(-x) = \frac{-ax+15}{-x-a} = \frac{ax-15}{x+a}$$

Điều kiện $f(-x) = -f(x)$ trở thành:

$$\frac{ax-15}{x+a} = -\frac{ax+15}{x-a} = \frac{-ax-15}{x-a}$$

Nhân chéo và rút gọn:

$$(ax-15)(x-a) = (-ax-15)(x+a)$$

$$ax^2 - a^2x - 15x + 15a = -ax^2 - a^2x - 15x - 15a$$

$$2ax^2 + 30a = 0 \implies 2a(x^2+15) = 0$$

Vì $x^2 + 15 > 0$ với mọi $x$ nên $a = 0$.

Khi đó $f(x) = \dfrac{15}{x}$, đúng là hàm lẻ.

\textbf{Chú ý.} Với $a = \pm\sqrt{15}$, mẫu số có nghiệm trùng với nghiệm của tử, hàm không xác định tại đó và không thỏa mãn định nghĩa hàm lẻ.

\textbf{Chú ý về kho đề.} Câu gốc dùng ở vị trí này (\texttt{CAE-M-CD4-08.1-E-003}) có đáp án không chặt chẽ: tập xác định đúng của $\dfrac{1}{\sqrt{x}-2}$ là $[0,4) \cup (4,+\infty)$, không đáp án nào trong bốn lựa chọn mô tả đúng. Câu đã được thay bằng \texttt{CAE-M-CD4-08.8-E-098}.
\end{dapan}

\begin{dapan}[21 — Đáp án D]
\textbf{Giải.} Hàm $f(x) = \dfrac{x^2-1}{3x-1}$ chỉ cần mẫu khác $0$:

$$3x - 1 \neq 0 \implies x \neq \frac{1}{3}$$

Tập xác định:

$$\left(-\infty, \frac{1}{3}\right) \cup \left(\frac{1}{3}, +\infty\right)$$

\textbf{Chú ý.} Đáp án A và B giống hệt nhau (đều ghi $\mathbb{R}$). Khi gặp hai đáp án trùng nhau trong đề trắc nghiệm, chắc chắn cả hai đều sai.
\end{dapan}

\begin{dapan}[22 — Đáp án B]
\textbf{Giải.} Hàm $y = f(x)$ là hàm lẻ nên $f(-a) = -f(a)$ với mọi $a$.

Suy ra điểm $(-a, -f(a))$ thỏa mãn $y = f(x)$ tại $x = -a$: tung độ $-f(a)$ đúng bằng $f(-a)$.

\textbf{Cách nhanh.} Điểm trên đồ thị hàm lẻ có dạng $(x, f(x))$. Thay $x = -a$ được điểm $(-a, f(-a)) = (-a, -f(a))$.
\end{dapan}

\begin{dapan}[23 — Đáp án C]
\textbf{Giải.} Cần tìm $x$ sao cho $\dfrac{x}{2x-1} = \dfrac{1}{3}$:

$$3x = 2x - 1 \implies x = -1$$

Thay $x = -1$ vào biểu thức:

$$f\!\left(\frac{1}{3}\right) = 2(-1)^2 + 3 = 2 + 3 = 5$$

\textbf{Bẫy.} Học sinh hay thay $x = \dfrac{1}{3}$ trực tiếp vào $2x^2+3$ và được $\dfrac{2}{9}+3$. Sai, vì $\dfrac{1}{3}$ là \textbf{giá trị của biểu thức trong ngoặc}, không phải giá trị của $x$.
\end{dapan}

\subsection*{Nhóm D — Ôn tích lũy}

\begin{dapan}[24 — Đáp án D]
\textbf{Giải.} Tập $M$ gồm các số có dạng $3k+1$ với $k \in \mathbb{Z}$. Kiểm tra từng đáp án bằng cách xét số dư khi chia cho $3$:
\begin{itemize}[leftmargin=1.4em, itemsep=2pt]
\item $0 = 3(0) + 0$: số dư $0$ ✗
\item $2 = 3(0) + 2$: số dư $2$ ✗
\item $3 = 3(1) + 0$: số dư $0$ ✗
\item $4 = 3(1) + 1$: số dư $1$ ✓
\end{itemize}


Vậy $4 \in M$.

\textbf{Cách nhanh.} Chỉ cần tính số dư của từng đáp án khi chia cho $3$, tìm đáp án có số dư bằng $1$.
\end{dapan}

\begin{dapan}[25 — Đáp án C]
\textbf{Giải.} Tính $a \odot 2 = a \cdot 2 - 2 = 2a - 2$ với từng $a \in A = \{0, 1, 2\}$:
\begin{itemize}[leftmargin=1.4em, itemsep=2pt]
\item $a = 0$: $0 - 2 = -2$
\item $a = 1$: $2 - 2 = 0$
\item $a = 2$: $4 - 2 = 2$
\end{itemize}


Vậy $B = \{-2, 0, 2\}$.

$$A \cap B = \{0,1,2\} \cap \{-2,0,2\} = \{0, 2\}$$

\textbf{Chú ý.} Đáp án B ($\{1\}$) là kết quả khi lấy nhầm $A \setminus B$ hoặc tính sai phép toán. Phải thay từng phần tử của $A$ rồi lấy giao.
\end{dapan}

\subsection*{Nhóm D — Ôn tích lũy}

\begin{dapan}[1 — Đáp án B]
\textbf{Giải.} Áp dụng công thức cộng của tan với $\tan\dfrac{\pi}{4} = 1$:

$$\tan\!\left(\alpha + \frac{\pi}{4}\right) = \frac{\tan\alpha + \tan\frac{\pi}{4}}{1 - \tan\alpha \cdot \tan\frac{\pi}{4}} = \frac{2+1}{1 - 2 \cdot 1} = \frac{3}{-1} = -3$$

\textbf{Bẫy.} Học sinh hay quên dấu trừ ở mẫu và chọn đáp án A. Công thức cộng của tan có mẫu là $1 - \tan\alpha\tan\beta$ (dấu ngược với tử).
\end{dapan}

\begin{dapan}[2 — Đáp án B]
\textbf{Giải.} Kiểm tra tính chẵn của từng hàm:
\begin{itemize}[leftmargin=1.4em, itemsep=2pt]
\item $y = x+1$: $f(-x) = -x+1 \neq f(x)$, không chẵn không lẻ ✗
\item $y = x^2+1$: $f(-x) = x^2+1 = f(x)$, hàm chẵn ✓
\item $y = \sin x$: hàm lẻ ✗
\end{itemize}

- $y = 2x$: hàm lẻ ✗
\end{dapan}

\begin{dapan}[3 — Đáp án C]
\textbf{Giải.} Hàm $y = \tan x = \dfrac{\sin x}{\cos x}$ xác định khi $\cos x \neq 0$:

$$x \neq \frac{\pi}{2} + k\pi, \quad k \in \mathbb{Z}$$

\textbf{Bẫy.} Đáp án B loại các điểm $x = k\pi$, nhưng đó là các điểm cos \textbf{bằng} $\pm 1$ (tan xác định). Đáp án D chỉ loại nửa số điểm vì dùng chu kỳ $2\pi$ thay vì $\pi$.
\end{dapan}

\begin{dapan}[4 — Đáp án C]
\textbf{Giải.} Áp dụng công thức cung phụ với $\dfrac{\pi}{2} - \alpha$ (ở đây $\alpha = x - \dfrac{\pi}{2}$ không áp dụng trực tiếp, dùng công thức cộng):

$$\cos\!\left(x - \frac{\pi}{2}\right) = \cos x\cos\frac{\pi}{2} + \sin x\sin\frac{\pi}{2} = \cos x \cdot 0 + \sin x \cdot 1 = \sin x$$

\textbf{Cách nhanh.} Trừ $\dfrac{\pi}{2}$ trong hàm cos thì đổi thành sin. Quy tắc: cos lệch $\dfrac{\pi}{2}$ thành sin, dấu phụ thuộc chiều dịch chuyển.
\end{dapan}

\begin{dapan}[5 — Đáp án B]
\textbf{Giải.} Cần tìm $\alpha$ sao cho $\tan\alpha = \sqrt{3}$. Giá trị này tương ứng với góc đặc biệt:

$$\tan\frac{\pi}{3} = \sqrt{3}$$

Vậy $\alpha = \dfrac{\pi}{3}$ là một giá trị thỏa mãn.

\textbf{Chú ý.} Các đáp án khác: $\tan\dfrac{\pi}{6} = \dfrac{\sqrt{3}}{3}$, $\tan\dfrac{\pi}{4} = 1$, $\tan\dfrac{2\pi}{3} = -\sqrt{3}$.
\end{dapan}

\begin{dapan}[6 — Đáp án D]
\textbf{Giải.} Viết $210^\circ = 180^\circ + 30^\circ$ và dùng cung hơn kém $\pi$:

$$\cos 210^\circ = \cos(180^\circ + 30^\circ) = -\cos 30^\circ = -\frac{\sqrt{3}}{2}$$

\textbf{Cách nhanh.} $210^\circ$ thuộc phần tư thứ ba, ở đó cos âm. Giá trị tuyệt đối bằng $\cos 30^\circ = \dfrac{\sqrt{3}}{2}$.
\end{dapan}

\begin{dapan}[7 — Đáp án A]
\textbf{Giải.} Khai triển hai công thức cộng:

$$\sin(\alpha+\beta) = \sin\alpha\cos\beta + \cos\alpha\sin\beta = \frac{3}{5}$$

$$\sin(\alpha-\beta) = \sin\alpha\cos\beta - \cos\alpha\sin\beta = \frac{1}{5}$$

Cộng hai phương trình:

$$2\sin\alpha\cos\beta = \frac{4}{5} \implies \sin\alpha\cos\beta = \frac{2}{5}$$

Trừ hai phương trình:

$$2\cos\alpha\sin\beta = \frac{2}{5} \implies \cos\alpha\sin\beta = \frac{1}{5}$$

Lập tỉ số:

$$\tan\alpha\cot\beta = \frac{\sin\alpha\cos\beta}{\cos\alpha\sin\beta} = \frac{2/5}{1/5} = 2$$

\textbf{Cách nhanh.} Cộng và trừ hai phương trình để tách riêng hai tích, không cần giải hệ tìm từng góc.
\end{dapan}

\begin{dapan}[8 — Đáp án A]
\textbf{Giải.} Chia cả tử và mẫu cho $\cos\alpha$:

$$\frac{2\sin\alpha + \cos\alpha}{\sin\alpha + 2\cos\alpha} = \frac{2\tan\alpha + 1}{\tan\alpha + 2} = \frac{2(7)+1}{7+2} = \frac{15}{9} = \frac{5}{3}$$

\textbf{Cách nhanh.} Cùng dạng với Câu 1 nhưng hệ số khác. Quy trình luôn là chia cho $\cos\alpha$ rồi thay $\tan\alpha$.
\end{dapan}

\begin{dapan}[9 — Đáp án D]
\textbf{Giải.} Với mọi $\alpha$, ta có $|\sin\alpha| \le 1$, tức $-1 \le \sin\alpha \le 1$.

Kiểm tra đáp án D: $\dfrac{\sqrt{5}}{2} \approx 1{,}118 > 1$. Vậy giá trị này không thể là $\sin\alpha$.

Ba đáp án còn lại đều nằm trong $[-1,1]$ nên đều có thể.

\textbf{Cách nhanh.} Bình phương từng đáp án và so với $1$: chỉ $\dfrac{5}{4} > 1$.
\end{dapan}

\begin{dapan}[10 — Đáp án D]
\textbf{Giải.} Xét từng mệnh đề:
\begin{itemize}[leftmargin=1.4em, itemsep=2pt]
\item A đúng: độ và radian là hai đơn vị đo góc.
\item B đúng: $1^\circ = \dfrac{360^\circ}{360}$, và $1$ rad $= \dfrac{2\pi \text{ rad}}{2\pi}$.
\item C đúng: $1$ rad $\approx 57{,}3^\circ > 1^\circ$.
\end{itemize}

- D sai: số đo radian của một góc \textbf{không phụ thuộc} bán kính. Công thức $l = |\alpha|r$ cho thấy tỉ số $\dfrac{l}{r}$ không đổi khi $r$ thay đổi.
\end{dapan}

\begin{dapan}[11 — Đáp án A]
\textbf{Giải.} Hàm $f(x) = \sin x + a\cos x$ có giá trị lớn nhất là $\sqrt{1+a^2}$. Đề cho giá trị này bằng $\sqrt{5}$:

$$\sqrt{1+a^2} = \sqrt{5} \implies 1+a^2 = 5 \implies a^2 = 4$$

Vì $a > 0$ nên $a = 2$. Hàm trở thành $f(x) = \sin x + 2\cos x$, đạo hàm $f'(x) = \cos x - 2\sin x$.

Điều kiện $f'(x_0) = -\dfrac{2}{5}$:

$$\cos x\_0 - 2\sin x\_0 = -\frac{2}{5}$$

Đặt $s = \sin x_0$, suy ra $\cos x_0 = 2s - \dfrac{2}{5}$. Thay vào $\sin^2 x_0 + \cos^2 x_0 = 1$:

$$s^2 + \left(2s - \frac{2}{5}\right)^2 = 1 \implies 5s^2 - \frac{8}{5}s - \frac{21}{25} = 0$$

Giải ra hai nghiệm: $s = \dfrac{3}{5}$ và $s = -\dfrac{7}{25}$.

\begin{itemize}[leftmargin=1.4em, itemsep=2pt]
\item Với $s = \dfrac{3}{5}$: $\cos x_0 = \dfrac{4}{5}$, $\tan x_0 = \dfrac{3}{4}$.
\item Với $s = -\dfrac{7}{25}$: $\cos x_0 = -\dfrac{24}{25}$, $\tan x_0 = \dfrac{7}{24}$.
\end{itemize}


Trong hai giá trị này, chỉ $\dfrac{3}{4}$ có mặt trong bốn lựa chọn. Vậy đáp án là A.

\textbf{Chú ý.} Kiểm tra lại nghiệm thứ nhất: $\cos x_0 - 2\sin x_0 = \dfrac{4}{5} - \dfrac{6}{5} = -\dfrac{2}{5}$, đúng bằng dữ kiện đề cho. Nghiệm thứ hai cũng thỏa mãn phương trình nhưng cho $\tan x_0 = \dfrac{7}{24}$, không nằm trong đáp án, nên đề chỉ nhắm vào nghiệm thứ nhất.
\end{dapan}

\begin{dapan}[12 — Đáp án C]
\textbf{Giải.} Hàm $y = \dfrac{1}{3}\sin 3x$ có dạng $\sin(\omega x)$ với $\omega = 3$. Chu kỳ:

$$T = \frac{2\pi}{3}$$

Tần số là nghịch đảo của chu kỳ:

$$f = \frac{1}{T} = \frac{3}{2\pi}$$

\textbf{Chú ý.} Biên độ $\dfrac{1}{3}$ không ảnh hưởng đến tần số. Chỉ hệ số của $x$ quyết định.
\end{dapan}

\begin{dapan}[13 — Đáp án A]
\textbf{Giải.} Đạo hàm của $y = \cos x$ là $y' = -\sin x$. Hệ số góc của tiếp tuyến tại $P\left(\dfrac{\pi}{3}, \dfrac{1}{2}\right)$:

$$y'\!\left(\frac{\pi}{3}\right) = -\sin\frac{\pi}{3} = -\frac{\sqrt{3}}{2}$$

Đường thẳng cần tìm \textbf{vuông góc} với tiếp tuyến nên hệ số góc của nó là nghịch đảo đối dấu:

$$k = \frac{1}{\sqrt{3}/2} = \frac{2}{\sqrt{3}} = \frac{2\sqrt{3}}{3}$$

Phương trình đường thẳng qua $P$:

$$y - \frac{1}{2} = \frac{2\sqrt{3}}{3}\left(x - \frac{\pi}{3}\right)$$

Biến đổi về dạng tổng quát, nhân cả hai vế với $\sqrt{3}$:

$$\sqrt{3}y - \frac{\sqrt{3}}{2} = 2\left(x - \frac{\pi}{3}\right) = 2x - \frac{2\pi}{3}$$

$$2x - \sqrt{3}y - \frac{2\pi}{3} + \frac{\sqrt{3}}{2} = 0$$

\textbf{Chú ý.} Đáp án C sai dấu ở số hạng $\dfrac{\sqrt{3}}{2}$. Phải kiểm tra lại bằng cách thay tọa độ điểm $P$ vào phương trình.
\end{dapan}

\begin{dapan}[14 — Đáp án B]
\textbf{Giải.} Áp dụng công thức cung bù với $180^\circ = \pi$:

$$\sin(180^\circ - \alpha) = \sin\alpha = \frac{1}{2}$$

\textbf{Cách nhanh.} Cung bù giữ nguyên giá trị sin. Đây là hệ quả trực tiếp của $\sin(\pi - \alpha) = \sin\alpha$.
\end{dapan}

\begin{dapan}[15 — Đáp án C]
\textbf{Giải.} Gọi $R$ là bán kính đường tròn, góc ở tâm là $\alpha = 2$ rad, dây cung $AB = 2$.

Hạ đường cao $OH$ từ tâm $O$ xuống dây $AB$. Khi đó $H$ là trung điểm $AB$ nên $AH = 1$, và $\angle AOH = \dfrac{\alpha}{2} = 1$ rad.

Trong tam giác vuông $OHA$:

$$\sin 1 = \frac{AH}{OA} = \frac{1}{R} \implies R = \frac{1}{\sin 1}$$

Độ dài cung:

$$l = \alpha R = 2 \cdot \frac{1}{\sin 1} = \frac{2}{\sin 1}$$

\textbf{Chú ý.} Đáp án D ($2\sin 1$) là kết quả khi nhân thay vì chia. Đáp án B ($\sin 2$) là kết quả khi quên nhân với bán kính.
\end{dapan}

\clearpage
\sectiondautrang{Tổng kết lỗi thường gặp}
\addcontentsline{toc}{section}{Tổng kết lỗi thường gặp}

Bảng chẩn đoán dùng khi chữa bài và khi soạn đề kiểm tra. Cột ``Câu liên quan'' trỏ tới đề bài tập về nhà.

\bang{
\begin{tabular}{@{}p{5.6cm}p{4.2cm}p{5.2cm}@{}}
\toprule
\textbf{Lỗi} & \textbf{Câu liên quan} & \textbf{Cách cảnh báo học sinh} \\
\midrule
Quên dấu theo phần tư & Câu 3, 11 phần bắt buộc & Mỗi lần lấy căn từ bình phương, hỏi ngay: "Góc ở phần tư thứ mấy?" \\
Nhầm dấu công thức cộng của cos & Câu 12 phần bắt buộc & Viết to công thức, khoanh tròn dấu: "cos cộng thì trong ngoặc trừ" \\
Nhầm sin với cos trong công thức cộng & Câu 12 phần bắt buộc & Quy tắc: "sin thì sin trước, cos thì cos trước" \\
Quên nhân đạo hàm hàm hợp & Câu 6 phần bắt buộc & Đạo hàm $\sin u$ phải nhân thêm $u'$ \\
Lẫn lãi đơn với lãi kép & Câu 17 phần bắt buộc & "Phần trăm của giá trị hiện có" là lũy thừa, không phải nhân tuyến tính \\
Nhầm $\tan 2\alpha$ với $\cot 2\alpha$ & Câu 10 phần bắt buộc & Đọc kỹ tên hàm trong đề trước khi chọn công thức \\
Thay sai biến vào hàm hợp & Câu 23 phần bắt buộc & Giải phương trình trong ngoặc để tìm $x$ trước, rồi mới thay \\
Quên điều kiện cơ số logarit & Câu 14 phần bắt buộc & Cơ số phải dương và khác $1$, loại nghiệm âm \\
Nhầm cung phụ với cung bù & Câu 14 phần tự chọn & "Thấy $\dfrac{\pi}{2}$ thì đổi tên hàm, thấy $\pi$ thì giữ nguyên" \\
Quên rằng $|\sin\alpha| \le 1$ & Câu 9 phần tự chọn & Thấy giá trị lớn hơn $1$ là loại ngay \\
Nhầm chu kỳ với tần số & Câu 12 phần tự chọn & Tần số là nghịch đảo của chu kỳ, không phải chính chu kỳ \\
Quên đổi dấu khi lấy hệ số góc vuông góc & Câu 13 phần tự chọn & Đường vuông góc có hệ số góc là nghịch đảo đối dấu \\
Quên điều kiện hàm lẻ phải xác định trên tập đối xứng & Câu 20 phần bắt buộc & Hàm lẻ phải xác định tại cả $x$ và $-x$ \\
\bottomrule
\end{tabular}}

% Không dùng \vfill: nó đẩy chữ xuống sát đáy trang,
% nơi banner letterhead chiếm chỗ (281-297mm) và sẽ che mất chữ.
\vspace{0.8cm}
\begin{center}
{\footnotesize\color{cae_blue} Tài liệu nội bộ, CAE SHANGHAI}
\end{center}

\end{document}
